\documentclass[12pt,letterpaper]{article}

% --- Paquetes de configuración regional y tipográfica ---
\usepackage[utf8]{inputenc}
\usepackage[spanish,es-tabla]{babel}
\usepackage[margin=2.5cm]{geometry}
\usepackage{amsmath,amsfonts,amssymb}
\usepackage{booktabs}
\usepackage{array}
\usepackage{xcolor}
\usepackage{listings}
\usepackage{caption}
\usepackage{fancyhdr}
\usepackage{hyperref}
\usepackage{float}
\usepackage{microtype}
\usepackage{tcolorbox}

% --- Colores institucionales y temáticos ---
\definecolor{codegreen}{rgb}{0.13,0.55,0.13}
\definecolor{codegray}{rgb}{0.45,0.45,0.45}
\definecolor{uctblue}{rgb}{0.08,0.22,0.45}
\definecolor{boxbg}{rgb}{0.96,0.97,0.99}

% --- Cajas destacadas ---
\tcbset{
    estudiobox/.style={
        colback=boxbg,
        colframe=uctblue,
        coltitle=white,
        fonttitle=\bfseries,
        boxrule=1pt,
        arc=3mm,
        left=4mm,
        right=4mm,
        top=3mm,
        bottom=3mm
    }
}

% --- Configuración de encabezados y pies de página ---
\pagestyle{fancy}
\fancyhf{}
\fancyhead[L]{\small\color{gray}Sistemas Inteligentes (INFO1157) -- Proyecto \#1}
\fancyhead[R]{\small\color{gray}Simulador Inercial 3D WitMotion}
\fancyfoot[C]{\thepage}
\renewcommand{\headrulewidth}{0.5pt}

% --- Hipervínculos ---
\hypersetup{
    colorlinks=true,
    linkcolor=uctblue,
    citecolor=uctblue,
    urlcolor=uctblue
}

\begin{document}

% =========================================================================
% PORTADA
% =========================================================================
\begin{titlepage}
    \centering
    \vspace*{1.5cm}
    {\Large\textbf{UNIVERSIDAD CATÓLICA DE TEMUCO}}\\[0.3cm]
    {\large FACULTAD DE INGENIERÍA}\\[0.2cm]
    {\large DEPARTAMENTO DE INGENIERÍA CIVIL INFORMÁTICA}\\[2.5cm]
    
    {\LARGE\bfseries INFORME TÉCNICO}\\[0.4cm]
    {\Large\bfseries PROYECTO \#1: SIMULADOR 3D DE SENSOR INERCIAL WITMOTION}\\[0.5cm]
    {\large Decodificación Binaria de Bajo Nivel (\texttt{ctypes}), Cinemática de Euler $Z-Y-X$ y Renderizado Interactivo en \texttt{VPython}}\\[2.5cm]
    
    \textbf{Asignatura:} Sistemas Inteligentes (INFO1157)\\[0.2cm]
    \textbf{Docente:} Dr. Alberto Caro S.\\[0.2cm]
    \textbf{Unidad:} Unidad 2 -- Adquisición y Procesamiento de Señales Sensoriales\\[2.0cm]
    
    \begin{flushright}
        \textbf{Estudiante:} Ingeniería Civil Informática\\
        \textbf{Repositorio:} \url{https://github.com/Diegood16/Sistemas2}\\
        \textbf{Fecha:} Noviembre de 2026\\
        \textbf{Temuco, Chile}
    \end{flushright}
    \vfill
\end{titlepage}

\newpage
\tableofcontents
\newpage

% =========================================================================
% 1. RESUMEN Y OBJETIVOS
% =========================================================================
\section{Resumen Ejecutivo y Objetivos}

En robótica, drones y sistemas de navegación, la estimación de la orientación espacial en tiempo real es crucial. Los sensores inerciales inteligentes (IMU/AHRS) combinan acelerómetros, giróscopos y magnetómetros para calcular los ángulos de inclinación y giro en el espacio tridimensional mediante algoritmos de fusión sensorial.

El presente proyecto implementa una solución completa en Python para procesar un flujo de datos binario pregrabado (\texttt{inercial.bin}) capturado desde un sensor inercial \textbf{WitMotion BWT901CL}. El software realiza la decodificación binaria de bajo nivel mediante estructuras \texttt{ctypes}, valida la integridad de cada paquete con sumas de comprobación (\textit{Checksum}), convierte las señales digitales en magnitudes físicas sexagesimales ($\pm 180^\circ$) y reproduce el movimiento del sensor en un entorno tridimensional interactivo con \texttt{VPython}.

\subsection{Objetivos Principales}
\begin{enumerate}
    \item \textbf{Decodificación de Bajo Nivel:} Desempaquetar bloques contiguos de 11 bytes utilizando la clase \texttt{Structure} de \texttt{ctypes} de Python.
    \item \textbf{Validación de Integridad:} Implementar el algoritmo de Checksum acumulado módulo 256 para filtrar tramas corruptas.
    \item \textbf{Conversión Física:} Transformar enteros con signo de 16 bits en complemento a dos (\textit{Little-Endian}) a grados sexagesimales continuos.
    \item \textbf{Visualización Tridimensional:} Modelar la geometría del sensor (caja roja y ejes ortogonales) y orientarla en el espacio 3D según la secuencia intrínseca de Euler $Z-Y-X$ sin acumular deriva numérica (\textit{drift}).
    \item \textbf{Control Interactivo:} Desarrollar un sistema de navegación temporal bidireccional (pausa, paso a paso adelante y atrás, control de velocidad 1.0x/2.0x y reinicio) operable simultáneamente por teclado y botones en pantalla.
\end{enumerate}

% =========================================================================
% 2. FUNDAMENTOS CONCEPTUALES DEL SENSOR Y LA TRAMA
% =========================================================================
\section{Fundamentos Conceptuales: Sensor y Trama Binaria}

\subsection{Características del Sensor WitMotion BWT901CL}
Es un sensor inercial inteligente de 9 grados de libertad (9-DOF). En su interior contiene:
\begin{itemize}
    \item Un microprocesador ARM Cortex-M0 que ejecuta internamente un Filtro de Kalman Extendido (EKF).
    \item Un acelerómetro triaxial (mide gravedad y aceleración lineal en $X, Y, Z$).
    \item Un giróscopo triaxial (mide velocidad angular en rad/s).
    \item Un magnetómetro triaxial (mide campo magnético terrestre para corregir el rumbo).
\end{itemize}

El microprocesador fusiona estas lecturas y entrega directamente los ángulos de actitud espacial: \textbf{Roll} (balanceo), \textbf{Pitch} (cabeceo) y \textbf{Yaw} (guiñada/rumbo).

\begin{tcolorbox}[estudiobox, title=Principio de Fusión y Salida Sensorial]
\textbf{Principio de Funcionamiento:} En lugar de transmitir aceleraciones crudas sujetas a ruido de alta frecuencia, el procesador interno del sensor calcula la orientación física exacta en el espacio tridimensional y emite paquetes binarios compactos de 11 bytes a través de Bluetooth o puerto serial a 10 Hz (10 veces por segundo).
\end{tcolorbox}

\subsection{Estructura Byte a Byte de la Trama de Ángulos (0x53)}
El archivo \texttt{inercial.bin} contiene una secuencia contigua de paquetes binarios de 11 bytes con la estructura detallada en la Tabla \ref{tab:trama}.

\begin{table}[H]
\centering
\caption{Estructura byte a byte de la trama de actitud angular (0x53).}
\label{tab:trama}
\begin{tabular}{cclp{7.5cm}}
\toprule
\textbf{Índice} & \textbf{Campo} & \textbf{Valor} & \textbf{Significado y Función} \\
\midrule
0 & \texttt{B0} & \texttt{0x55} (85) & Byte de cabecera fijo (\textit{Header}) para sincronizar el inicio del paquete. \\
1 & \texttt{B1} & \texttt{0x53} (83) & Identificador de tipo: indica que este paquete contiene Ángulos de Euler. \\
2 & \texttt{B2} & \texttt{Roll\_L} & Byte bajo (LSB) del ángulo de balanceo. \\
3 & \texttt{B3} & \texttt{Roll\_H} & Byte alto (MSB) del ángulo de balanceo. \\
4 & \texttt{B4} & \texttt{Pitch\_L} & Byte bajo (LSB) del ángulo de cabeceo. \\
5 & \texttt{B5} & \texttt{Pitch\_H} & Byte alto (MSB) del ángulo de cabeceo. \\
6 & \texttt{B6} & \texttt{Yaw\_L} & Byte bajo (LSB) del ángulo de guiñada. \\
7 & \texttt{B7} & \texttt{Yaw\_H} & Byte alto (MSB) del ángulo de guiñada. \\
8 & \texttt{B8} & \texttt{Temp\_L} & Byte bajo (LSB) de la temperatura del sensor. \\
9 & \texttt{B9} & \texttt{Temp\_H} & Byte alto (MSB) de la temperatura del sensor. \\
10 & \texttt{BA} & \texttt{Checksum} & Byte de control para verificar que ningún bit se corrompió. \\
\bottomrule
\end{tabular}
\end{table}

\subsection{El Algoritmo de Verificación (Checksum)}
Para garantizar que ningún byte se haya dañado en la transmisión por Bluetooth, el paquete incluye el byte \texttt{BA}:
\begin{equation}
\text{Checksum}_{\text{calculado}} = \left( B_0 + B_1 + B_2 + \dots + B_9 \right) \pmod{256}
\end{equation}
Si $\text{Checksum}_{\text{calculado}} == B_{10}$, el paquete es íntegro y legítimo. En el análisis de \texttt{inercial.bin}, las \textbf{406 tramas} pasaron exitosamente esta prueba ($100\%$ de validez).

\subsection{Fórmula de Conversión a Grados Sexagesimales}
Cada ángulo se almacena en 2 bytes (16 bits) bajo formato \textbf{Little-Endian} y complemento a dos con signo:
\begin{equation}
\text{ValorCrudo} = (\text{ByteAlto} \ll 8) \mid \text{ByteBajo} \quad \in [-32768, +32767]
\end{equation}
El fabricante estipula que este rango numérico digital de 16 bits cubre el intervalo angular completo de $\pm 180^\circ$:
\begin{equation}
\text{Ángulo} [^\circ] = \frac{\text{ValorCrudo}}{32768.0} \times 180.0
\end{equation}

% =========================================================================
% 3. ARQUITECTURA DEL SOFTWARE Y ENTORNO 3D
% =========================================================================
\section{Arquitectura del Software y Entorno Tridimensional}

El código se organizó en tres módulos limpios, manteniendo principios de alta cohesión y nombres naturales y autoexplicativos:

\begin{table}[H]
\centering
\caption{Módulos de software del Proyecto \#1.}
\label{tab:archivos_p1}
\begin{tabular}{llp{7.5cm}}
\toprule
\textbf{Archivo} & \textbf{Ubicación} & \textbf{Responsabilidad Principal} \\
\midrule
\texttt{decoder.py} & \texttt{src/} & Define la estructura \texttt{eTrama} con \texttt{ctypes}, valida el Checksum y convierte los 11 bytes a grados físicos. \\
\texttt{visualizer.py} & \texttt{src/} & Construye la escena 3D en \texttt{VPython}, el modelo del sensor (caja roja y ejes), el HUD y los botones interactivos. \\
\texttt{main.py} & \texttt{src/} & Bucle de animación a 10 Hz, controlador de eventos de teclado bidireccional, telemetría por consola y plantilla Bluetooth COM9. \\
\bottomrule
\end{tabular}
\end{table}

\subsection{Modelado Geométrico}
Siguiendo las especificaciones exactas del proyecto:
\begin{itemize}
    \item \textbf{Cuerpo del Sensor:} Una caja roja de dimensiones $(100, 60, 160)$ unidades.
    \item \textbf{Eje X (Transversal):} Barra cilíndrica azul de dimensiones $(200, 2, 2)$.
    \item \textbf{Eje Y (Vertical):} Barra cilíndrica verde de dimensiones $(2, 200, 2)$.
    \item \textbf{Plano Base de Referencia:} Un prisma azul oscuro debajo del sensor de dimensiones $(300, 10, 300)$.
\end{itemize}

\subsection{Rotación en Secuencia $Z-Y-X$ sin Deriva (\textit{Drift})}
Si en cada cuadro de la animación se rotara sobre la orientación anterior, los pequeños errores de punto flotante de la computadora se acumularían, haciendo que el sensor se deforme o se desvíe con el tiempo (\textit{drift}). 

Para solucionar esto, la función \texttt{orientar\_absoluto()} resetea el sensor a sus vectores base originales en cada fotograma y aplica las tres rotaciones de Tait-Bryan en orden estricto:
\begin{enumerate}
    \item Rotación de guiñada (\textbf{Yaw}, $\psi$) alrededor del eje $Z$.
    \item Rotación de cabeceo (\textbf{Pitch}, $\theta$) alrededor del eje $Y$.
    \item Rotación de balanceo (\textbf{Roll}, $\phi$) alrededor del eje $X$.
\end{enumerate}

% =========================================================================
% 4. RESULTADOS Y TELEMETRÍA EXPERIMENTAL
% =========================================================================
\section{Resultados y Análisis Experimental}

Al procesar las 406 tramas de \texttt{inercial.bin}, se obtuvo la serie temporal completa de los tres ángulos de orientación:
\begin{itemize}
    \item \textbf{Roll:} Muestra oscilaciones periódicas entre $-20^\circ$ y $+25^\circ$, correspondientes al balanceo transversal del sensor.
    \item \textbf{Pitch:} Registra inclinaciones longitudinales suaves confinadas dentro del rango $\pm 15^\circ$.
    \item \textbf{Yaw:} Presenta un giro progresivo continuo acumulando un cambio de rumbo de $180^\circ$.
    \item \textbf{Duración total:} $406 \text{ tramas} \div 10 \text{ Hz} = 40.6 \text{ segundos}$, coincidiendo exactamente con la duración de la grabación de referencia \texttt{MyDemo.mp4}.
\end{itemize}

% =========================================================================
% 5. MANUAL DE CONTROLES INTERACTIVOS
% =========================================================================
\section{Manual de Controles Interactivos}

El simulador ofrece control híbrido (teclas rápidas y botones gráficos en la ventana 3D):

\begin{table}[H]
\centering
\caption{Tabla de controles interactivos del simulador.}
\label{tab:controles_p1}
\begin{tabular}{llp{7cm}}
\toprule
\textbf{Acción} & \textbf{Teclado} & \textbf{Botón en Pantalla 3D} \\
\midrule
Pausar / Continuar & Barra Espaciadora & \texttt{[⏯️ Pausa / Play]} \\
Retroceder 1 fotograma & Flecha Izquierda / Tecla \texttt{B} & \texttt{[⏪ Paso Atrás]} \\
Avanzar 1 fotograma & Flecha Derecha / Tecla \texttt{T} & \texttt{[⏩ Paso Adelante]} \\
Velocidad Normal (1.0x) & Tecla \texttt{1} & \texttt{[⚡ 1.0x]} \\
Velocidad Rápida (2.0x) & Tecla \texttt{2} & \texttt{[⚡ 2.0x]} \\
Reiniciar animación & Tecla \texttt{R} & \texttt{[🔄 Reiniciar]} \\
\bottomrule
\end{tabular}
\end{table}

% =========================================================================
% 6. CONCLUSIONES
% =========================================================================
\section{Conclusiones}

El Proyecto \#1 permitió adquirir competencias prácticas fundamentales en la captura, decodificación binaria de bajo nivel y visualización tridimensional de sensores inerciales modernos. 

La separación arquitectónica entre decodificación binaria, cinemática de Euler y renderizado en \texttt{VPython} permitió lograr una simulación fluida a 10 Hz con $100\%$ de paquetes válidos. El agregado de controles temporales bidireccionales y la eliminación del drift mediante rotaciones absolutas $Z-Y-X$ aseguran una solución de ingeniería robusta, profesional y de fácil comprensión.

\end{document}
